\inicializarSeccion{Intervalos de Confianza}

Como se logró observar por medio de las pruebas de independencia y de hipótesis, se puede estimar el nivel de confianza para nuestros parámetros de interés. Utilizando los intervalos de confianza, podemos verificar si los resultados obtenidos son estadísticamente significativos y reiterar la validez de nuestras conclusiones.

Para la hipótesis de que el nivel promedio de oxígeno disuelto es mayor al umbral mínimo óptimo, podemos construir un intervalo de confianza para estimar el rango dentro del cual se encuentra el verdadero valor del promedio con un cierto nivel de confianza.

\begin{figure}[H]
    \centering
    \includegraphics[width=\linewidth]{parte2/intervaloDeConfianza1.png}
    \caption{Intervalo de Confianza para el Nivel Promedio de Oxígeno Disuelto}
    \label{fig:intervaloDeConfianza1}
\end{figure}


Como podemos observar en \autoref{fig:intervaloDeConfianza1}, debido a la gran cantidad de datos disponibles, y debido a que la desviación estandar de los datos es extremadamente bajo, el intervalo de confianza es muy estrecho, lo que indica una alta precisión en la estimación del promedio.

Analizando el segundo intervalo de confianza, podemos observar la relación entre el oxígeno disuelto y la densidad del agua. Este intervalo también es bastante estrecho, lo que sugiere una alta precisión en la estimación de esta relación que se muestra en \autoref{fig:intervaloDeConfianza2}.

\begin{figure}[h]
    \centering
    \includegraphics[width=\linewidth]{parte2/intervaloDeConfianza2.png}
    \caption{Intervalo de Confianza entre oxigeno disuelto y densidad del agua}
    \label{fig:intervaloDeConfianza2}
\end{figure}

Estos resultados no nos deberían de sorprender debido a la naturaleza de los datos y la baja variabilidad observada. La poca estrechez del intervalo coincide con la distancia de desviación estándar de los datos. Estas observaciones refuerzan la confianza en la precisión de nuestros estimadores.